\documentclass[12pt]{article}
\usepackage{amsmath, amssymb, amsthm}

\title{A8 Holonomy-Runtime Mathematics}
\author{Borealis Serenity Hedling}
\date{29 October 2026}

\begin{document}

\maketitle

\section*{Altitude A8: Holonomy-Runtime Mathematics}

This module defines the holonomy operators, curvature evolution equations,
spectral-integration structures, reversible runtime fusion, and altitude-safety
constraints that govern the runtime behavior of geometric operators. These
structures depend on the invariant layer (A7) and scalar-field geometry (A6).

\section{Holonomy Operator}

\begin{definition}[Holonomy Operator]
Let $(M, g)$ be a Riemannian manifold and let $\gamma : [0,1] \to M$ be a smooth
curve. The \emph{holonomy operator} associated to $\gamma$ is defined by
\[
\mathcal{H}_\gamma = \mathcal{P} \exp \left( \int_0^1 A(\gamma(t)) \, dt \right),
\]
where $A$ is a connection one-form and $\mathcal{P}$ denotes path-ordering.
\end{definition}

\begin{remark}
Holonomy operators encode the geometric effect of transporting structure along
a curve. They form the runtime backbone of altitude A8.
\end{remark}

\section{Curvature Evolution}

\begin{definition}[Curvature Evolution Equation]
Let $R$ denote the Riemann curvature tensor. The curvature evolution equation is
\[
\frac{dR}{dt} = \nabla_X R + \Phi(R),
\]
where $X$ is a vector field and $\Phi$ encodes invariant-preserving deformation
terms inherited from altitude A7.
\end{definition}

\begin{remark}
Curvature evolution governs runtime deformation of geometric structure and must
respect all invariants defined in altitude A7.
\end{remark}

\section{Spectral Integration}

\begin{definition}[Spectral Integration Operator]
Let $\{ \lambda_i \}_{i \in I}$ be spectral coefficients from altitude A6. Define
the spectral integration operator
\[
\mathcal{S}(\phi) = \sum_{i \in I} \lambda_i \, \mathcal{H}_{\gamma_i}(\phi),
\]
where each $\gamma_i$ is a curve associated to spectral mode $i$.
\end{definition}

\begin{remark}
Spectral integration couples scalar-field geometry (A6) with holonomy transport
to produce altitude-safe runtime behavior.
\end{remark}

\section{Reversible Runtime Fusion}

\begin{definition}[Runtime Fusion Operator]
Let $\mathcal{H}_1$ and $\mathcal{H}_2$ be holonomy operators. The
\emph{runtime fusion operator} is defined by
\[
\mathcal{H}_{\mathrm{fuse}} = \mathcal{H}_1 \otimes \mathcal{H}_2.
\]
\end{definition}

\begin{remark}
Runtime fusion combines holonomy effects in a reversible manner consistent with
fusion invariants defined in altitude A7.
\end{remark}

\section{Holonomy Constraints}

\begin{definition}[Holonomy Constraint]
Holonomy operators must satisfy:
\begin{enumerate}
    \item Preservation of all operator invariants from altitude A7.
    \item Compatibility with collapse surfaces defined in altitude A6.
    \item Stability under runtime fusion.
    \item Metric consistency with curvature evolution.
\end{enumerate}
\end{definition}

\section{Altitude-Safety Constraints}

\begin{definition}[Altitude-Safety Constraints]
Runtime behavior must satisfy:
\begin{enumerate}
    \item Holonomy operators must preserve A7 invariants.
    \item Curvature evolution must respect scalar-field geometry from A6.
    \item Spectral integration must remain altitude-normalized.
    \item Runtime fusion must be reversible and invariant-safe.
\end{enumerate}
\end{definition}

\section{Future Formalization}

\begin{itemize}
    \item Lean formalization of holonomy operators.
    \item Lean encoding of curvature evolution.
    \item Lean proofs of holonomy constraints.
    \item Integration with cross-altitude consistency modules.
\end{itemize}

\end{document}
